\exerciseeditionsection{Solutions to Wald's Problems}
\ExerciseEditorialNote

\begin{enumerate}[leftmargin=2.8em,itemsep=1.2em]
\WaldSolution{1}
Because \(\Omega\) is a coordinate,
\[
 \widetilde\nabla_\mu\widetilde\nabla_\nu\Omega
 =-\widetilde\Gamma^{\Omega}{}_{\mu\nu}.
\]
Thus every component \(\widetilde\Gamma^{\Omega}{}_{\mu\nu}\) vanishes at
\(\Omega=0\).  On \(\mathcal I^+\), the metric is
\[
 \dd\widetilde s^2=2\,\dd\Omega\,\dd u+
 \dd\theta^2+\sin^2\theta\,\dd\phi^2.
\]
Inserting these boundary values in the Christoffel formula gives, successively,
the vanishing of the first \(\Omega\)-derivatives of
\(\widetilde g_{uu}\), \(\widetilde g_{u\theta}\), and
\(\widetilde g_{u\phi}\).  Their values at \(\Omega=0\) also vanish.
Taylor expansion transverse to \(\mathcal I^+\) therefore yields
\[
 \widetilde g_{uu}=\order(\Omega^2),\qquad
 \widetilde g_{u\theta}=\order(\Omega^2),\qquad
 \widetilde g_{u\phi}=\order(\Omega^2).
\]

\WaldSolution{2}
Choose the stated normal coordinates at \(\Lambda\).  The \(C^{>0}\)
regularity of \(\widetilde h_{ij}\) gives
\[
 \widetilde h_{ij}=\delta_{ij}+\order(R),\qquad
 \partial_K\widetilde h_{ij}=\order(1).
\]
Condition (iii), with \(\Omega=0\), \(\widetilde D_i\Omega=0\), and the
normalized Hessian at \(\Lambda\), gives
\(\Omega=R^2[1+\order(R)]\).  Under inversion,
\(R=1/r\) and the Jacobian is \(\order(r^{-2})\).  Since
\(h_{ab}=\Omega^{-2}\widetilde h_{ab}\), the Jacobian factors cancel the
leading \(R^{-4}\), leaving
\[
 h_{ij}=\delta_{ij}+\order(r^{-1}),\qquad
 \partial_k h_{ij}=\order(r^{-2}).
\]
Condition (iv) gives one further inverse power for the physical extrinsic
curvature, so \(K_{ij}=\order(r^{-2})\).  The three-dimensional Ricci tensor
contains one derivative of the Christoffel symbols and quadratic Christoffel
terms.  Hence
\[
 {}^{(3)}R_{ij}=\order(r^{-3})+\order(r^{-4})
 =\order(r^{-3}).
\]

\WaldSolution{3}
\begin{enumerate}[label=\alph*),leftmargin=2em]
\item With \(u=t-r_*\), set \(\Omega=1/r\).  Multiplication of the
Schwarzschild metric by \(\Omega^2\) gives
\[
 \dd\widetilde s^2=-\Omega^2(1-2M\Omega)\dd u^2
 +2\,\dd u\,\dd\Omega+\dd\omega^2.
\]
This metric is smooth and nondegenerate at \(\Omega=0\), where
\(\widetilde\nabla_a\Omega\) is null.  The Schwarzschild solution is therefore
asymptotically flat at \(\mathcal I^+\).

\item A \(t=\text{constant}\) slice has
\[
 h=(1-2M/r)^{-1}\dd r^2+r^2\dd\omega^2,qquad K_{ij}=0.
\]
In inverted coordinates \(R=1/r\), \(\widetilde h=R^4h\) extends with
Euclidean value at \(R=0\), while \(\Omega=R^2\) has the required zero,
gradient, and Hessian there.  All conditions for an asymptotically flat initial
data set follow.
\end{enumerate}

\WaldSolution{4}
\begin{enumerate}[label=\alph*),leftmargin=2em]
\item Stationarity gives \(\partial_0T_{\mu\nu}=0\).  For spatial \(i\),
\[
 \partial_j(x^iT^{j}{}_{\nu})=T^i{}_{\nu}.
\]
The integral of the left side is a surface term and vanishes for localized
matter.  Thus \(\int T_{i\nu}\dd^3x=0\); symmetry also gives
\(\int T_{0i}\dd^3x=0\).  Only \(\int T_{00}\dd^3x\) can remain.

\item To first order, \(n^a=\xi^a=(\partial_t)^a\), \(\dd V=\dd^3x\),
and \(T=-T_{00}+\sum_iT_{ii}\).  Equation~\eqref{eq:11.2.10} becomes
\[
 M=2\int\left(T_{00}+\tfrac12T\right)\dd^3x
 =\int\left(T_{00}+\sum_iT_{ii}\right)\dd^3x
 =\int T_{00}\dd^3x.
\]
\end{enumerate}

\WaldSolution{5}
\begin{enumerate}[label=\alph*),leftmargin=2em]
\item For \(\xi^a=(\partial_t)^a\) in Schwarzschild,
\(\nabla^r\xi^t=M/r^2\).  Substitution in the Komar sphere integral gives
\[
 -\frac1{8\pi}\int_{S_r}\epsilon_{abcd}\nabla^c\xi^d=M,
\]
independently of \(r\) in the vacuum exterior.

\item For a perfect fluid, \(n^a=e^{-\phi}\xi^a\), and the Komar integrand
reduces to \((\rho+3p)e^\phi\).  Equation~\eqref{eq:6.2.10}, expressed with
the proper volume element
\(\dd V=4\pi r^2(1-2m/r)^{-1/2}\dd r\), gives the other expression
\[
 M=\int_\Sigma\rho(1-2m/r)^{1/2}\dd V.
\]
Equality produces the displayed relation in the problem.

\item Restore units and expand
\(e^{\phi/c^2}=1+\phi/c^2+\order(c^{-4})\) and
\((1-2Gm/rc^2)^{1/2}=1-Gm/(rc^2)+\order(c^{-4})\).
After the common rest-mass term cancels, the order \(c^{-2}\) equation is
\[
 3\int p\,\dd V+\frac12\int\rho\phi\,\dd V=0.
\]
The second term is the Newtonian gravitational energy \(E_g\), while
\(K=\tfrac32\int p\dd V\).  Therefore \(E_g=-2K\).
\end{enumerate}

\WaldSolution{6}
\begin{enumerate}[label=\alph*),leftmargin=2em]
\item The Killing identity gives
\(\nabla_a\nabla^a\psi^b=-R^b{}_c\psi^c\).  Stokes' theorem between two
vacuum spheres makes their surface integrals equal.  Applying it between an
inner boundary and infinity, then using Einstein's equation and
\(n_a\psi^a=0\), yields
\[
 J=-\int_\Sigma T_{ab}n^a\psi^b\dd V.
\]

\item A static slice is orthogonal to the timelike Killing field.  Staticity
removes the mixed time--azimuthal momentum density, so
\(T_{ab}n^a\psi^b=0\).  The volume formula gives \(J=0\).

\item At large \(r\), charged Kerr has
\(g_{t\phi}=-2Ma\sin^2\theta/r+\order(r^{-2})\).  Only this leading term
survives the angular-momentum surface integral.  Its sphere integral gives
\(J=Ma\); the charge-dependent terms decay one power faster.
\end{enumerate}
\end{enumerate}
